\lecture{11}{Fri 23 Sep 2022 11:34}{More on Connectedness}

Let $U \subset \mathbb{R}^p$ and  $a,b \in U$. We say that there is a \textit{polygonal curve} from $a$ to $b$ in $U$ if there are $z_0,z_1,\ldots,z_N \in U$ such that $z_0 = a$, $z_N = b$, and $[z_{i-1 }, z_i] \subset U$.

\begin{theorem}
	Open set $U \subset \mathbb{R}^p$ is connected if and only if for any $a, b \in U$, there is a polygonal curve from $a$ to $b$.
\end{theorem}
\begin{proof}
	Suppose $U$ is not connected. This means there is a disconnection $A,B$. Pick $a \in A \cap U$, and $b \in B \cap U$. By assumption, there is a polygonal curve from $a$ to $b $, $z_0,z_1,\ldots,z_N$.

	We have $z_0 \in A$ but $z_N \not\in A$. So there is a smallest index $n \ge 1$ such that $z_n \not\in A$, but $z_{n-1 } \in A$.

	Then  look at 
	\begin{align*}
		\varphi: \mathbb{R} &\longrightarrow \mathbb{R}^p \\
		t &\longmapsto \varphi(t) = (1-t)z_{n-1} + tz_n
	.\end{align*}
Let $A_1 = \{\varphi^{-1}(A) = \{t: \varphi(t) \in A\} \} $, and 
	$B_1 = \{\varphi^{-1}(B) = \{t: \varphi(t) \in B\} \} $. 
	Repeat the argument in Corollary \ref{thm: Rp-connected}. Thus 
	$A_1, B_1 $ form a disconnection of $[0,1]$.This is a contradiction since $[0,1]$ is connected.

	\begin{remark}
		In this part of the proof, we haven't used the openness. For example, every convex set is connected.
	\end{remark}

	Now prove the other direction. We say $a \sim b$ if there is a polygonal curve from  $a$ to $b$ in $U$. We check $\sim $ is a equivalence property:
	\begin{itemize}
		\item $a \sim a$. Indeed, $z_0 = a$ satisfies.
		\item $a\sim b \implies b \sim a$. Indeed, we reverse the order of $z_0,\ldots,z_N$.
		\item $a\sim b, b\sim c \implies a\sim c$. Indeed. We connect the curves.
	\end{itemize}
	Let $U_a = \{b \in U: b \sim a\} $.
	\begin{claim}
		$U_a$ is open.
	\end{claim}
	\begin{proof}
		Let $b \in U_a$. Let $z_0,\ldots,z_N$ be polygonal curve from $a$ to $b$. Since $U$ open, there exists $\delta>0$ such that $B_\delta(b) \subset U$. Then for every $c \in B_\delta(b)$ we have $c \sim b$, since $[c,b] \subset B_\delta(b)$ (since balls are convex). This means $a \sim c$, hence $B_\delta(b) \subset U_a$. Therefore  $U_a$ open.
	\end{proof}

	Let $V_a = \{b \in U: b \not\sim a\} $. Similarly,
	\begin{claim}
		$V_a$ is open.
	\end{claim}
	\begin{proof}
		Same as for $U_a$.
	\end{proof}

	Moreover, we see that $U \subset U_a \cup V_a$.
	If $U_a \cap V_a \neq \emptyset $, then both $U_a$, $V_a$ are noempty, they will form a disconnection, a contradiction. Hence $U = U_a$.
\end{proof}

\begin{corollary}
	Every open set is a union of disjoint connected open sets.
\end{corollary}

\begin{lemma}
	Disjoint family of open set can be at most countable.
\end{lemma}
\begin{proof}
	Use rational points. Now two open sets can contain the same rational ball.
\end{proof}


\begin{theorem}
	Every open set is a disjoint, countable union of connected sets.

	They are called \textit{connected components} of the open set.
\end{theorem}
\begin{example}
	Look at $F^c$, $F$ is the cantor set.
\end{example}

\section{Sequences}

\begin{definition}[Sequence]
	A sequence $E$ is a mapping \[
		\varphi: \mathbb{N} \to  E
	.\] 
	We identify $\varphi$ with ordered list \[
		(\varphi(1), \varphi(2), \ldots, \varphi(n), \ldots)
	.\] 
\end{definition}
\begin{notation}
	Often we use subscripts $x_n := \varphi(n)$.
	We may write $X = (x_1, x_2, \ldots)$, $(x_n)_{n\in \mathbb{N}}$, $(x_n)$, $(x_n)_{n = 1}^\infty $, or $\left\{ x_n \right\}_{n=1}^\infty  $ for the sequence. 
\end{notation}
